\section{Conclusion}
\label{sec:conclusion}

This paper has characterised the behaviour of machine learning detectors of GNSS spoofing
under attacks that are transmitted signals selected by an adversary. Each attack was
synthesised as radio samples, summed with authentic recordings from the Texas Spoofing Test
Battery, passed through a model of the receiver front end and tracked by an open software
receiver, so that the observables each detector read were produced by receiver processing
of a transmittable waveform. The generator was validated against a recorded attack whose
construction is documented in full, and reproduced its observables satellite by satellite.

Fourteen detectors, eight classical and six deep, were placed at a common operating point
defined by a recall floor of 0.95 on clean validation data, and evaluated against 126
transmitter settings satisfying the published conditions for receiver takeover.

Clean detection performance was found to be a poor predictor of resistance. The random
forest attained the highest clean $F_1$ score at 0.833 and the lowest false alarm rate at
0.337, and failed to flag 55 of the 126 attacks, flagging 0.1 percent of the spoofed
samples of one of them. Across the set, the four detectors with the highest clean $F_1$
were evaded at 0.226 of their detector-attack pairs against 0.106 for the remaining ten,
significant at $p = 2.8 \times 10^{-10}$, although the monotone rank correlation across all
fourteen is not significant and is reported as such. The
five detectors that flagged every attack did so while flagging between 46 and 74 percent of
authentic samples, a regime in which the resistance is not operationally available. Across
all fourteen detectors, false alarm rates at this operating point ranged from 0.337 to
0.960.

One evaded attack was propagated through the navigation stage. It drove the receiver clock
at 48.89 nanoseconds per second against a commanded 48.87, while the detector with the best
clean performance classified 98.8 percent of its samples as authentic, establishing
detection failure and navigation domain effect on the same attack and the same recording.

The adversary producing these results had no knowledge of the deployed detector, computed
no gradient and issued no query, and was further restricted to a single signal, full
satellite coverage and a stationary victim. Each restriction reduces its capability, so the
reported failure rates constitute a lower bound.

The broader implication is that detection accuracy measured on a fixed corpus of recorded
attacks characterises that corpus. A detector intended to withstand an adversary requires
characterisation against signals the adversary can select and emit, reported jointly with
the false alarm rate incurred. The library of transmitted attacks constructed here is
available for that purpose, and extends directly to the evaluation of physics constrained
and certified detectors.
